
\documentclass{article}
\usepackage{colortbl}
\usepackage{makecell}
\usepackage{multirow}
\usepackage{supertabular}

\begin{document}

\newcounter{utterance}

\centering \large Interaction Transcript for game `wordle\_withcritic', experiment `medium\_frequency\_words\_clue\_with\_critic', episode 14 with reflexion{-}llama{-}3.1{-}8b{-}instruct.
\vspace{24pt}

{ \footnotesize  \setcounter{utterance}{1}
\setlength{\tabcolsep}{0pt}
\begin{supertabular}{c@{$\;$}|p{.15\linewidth}@{}p{.15\linewidth}p{.15\linewidth}p{.15\linewidth}p{.15\linewidth}p{.15\linewidth}}
    \# & $\;$A & \multicolumn{4}{c}{Game Master} & $\;\:$B\\
    \hline

    \theutterance \stepcounter{utterance}  
    & & \multicolumn{4}{p{0.6\linewidth}}{
        \cellcolor[rgb]{0.9,0.9,0.9}{
            \makecell[{{p{\linewidth}}}]{
                \texttt{\tiny{[P1$\langle$GM]}}
                \texttt{You have attempted this same task before and did not succeed. These are the plans you wrote after those attempts.} \\
\texttt{Trial \#1: New plan: The last episode got aborted because my guess "work" is not a valid response. I should have produced a response in the format "explanation: details" followed by "guess: word". Next time, I will make sure to follow the correct format and provide a five{-}letter word related to work or daily activities, for example, "daily".} \\
\\ 
\texttt{Trial \#2: New plan: The last episode got aborted because my guess "office" is not a valid response. I should have produced a response in the format "explanation: details" followed by "guess: word". Moreover, in my last explanation I planned to find a new word related to work or daily activities, but I didn't consider the agreement with the previous guess "daily". Next time, I will make sure to follow the correct format and provide a five{-}letter word related to work or daily activities, and also take into account the agreement with the previous guess. For example, I could try to modify the guess "daily" to fit the clue, such as "duty".} \\
\\ 
\texttt{Trial \#3: New plan: The last episode got aborted because my guess "duty" is not a valid response. I should have produced a response in the format "explanation: details" followed by "guess: word". Moreover, in my last explanation I planned to modify the guess "daily" to fit the clue, but I didn't consider the fact that the previous guess was agreed upon. Next time, I will make sure to follow the correct format and provide a five{-}letter word related to work or daily activities, and also take into account the agreement with the previous guess. Since the previous guess "daily" was agreed upon, I will try to modify it by changing one letter at a time, for example, "dally".} \\
\\ 
\\ 
\texttt{You are a language wizard who likes to guess words by using the given rules.} \\
\\ 
\texttt{Welcome to Wordle! You have six attempts to guess the target word, a valid English word of five lowercase letters (a{-}z). Please use the tags "explanation:" and "guess:" to provide a concise explanation for each guess.} \\
\\ 
\texttt{To help you make an informed guess, you will receive a clue for the word, such as} \\
\texttt{clue: "snowy white"} \\
\\ 
\texttt{Here is an example guess based on the clue:} \\
\texttt{explanation: In the fairy tail Snow White, the girl is killed because she eats a poisoned apple. And the word apple has 5 letters.} \\
\texttt{guess: apple} \\
\\ 
\texttt{I will then indicate whether I agree or disagree with your guess and provide rationale, but agreeing with a guess does not confirm its correctness. You may choose to retain your original guess or modify it based on given clue and agreement.} \\
\\ 
\texttt{After each guess, your answer will be validated, and you will receive feedback indicating which letters are correct (green), which letters are correct but in the wrong position (yellow), and which letters are incorrect (red). This feedback can be useful in determining which letters to include or exclude in your next guess.} \\
\\ 
\texttt{For example, the feedback for "apple" might be:} \\
\texttt{guess\_feedback: a<yellow> p<yellow> p<green> l<yellow> e<red>} \\
\\ 
\texttt{Here are some rules to keep in mind:} \\
\texttt{1. Please do not generate any guess feedback or clue on your own} \\
\texttt{2. Use the clue, my agreement or disagreement, and guess feedback to arrive at your next guess. This should be reflected in your explanation as well.} \\
\\ 
\texttt{Do not generate any other text.} \\
\\ 
\texttt{Let's begin with your first guess.} \\
\\ 
\texttt{clue: nine{-}to{-}five routine} \\
            }
        }
    }
    & \\ \\

    \theutterance \stepcounter{utterance}  
    & \multicolumn{4}{p{0.6\linewidth}}{
        \cellcolor[rgb]{0.9,0.9,0.9}{
            \makecell[{{p{\linewidth}}}]{
                \texttt{\tiny{[P1$\rangle$GM]}}
                \texttt{explanation: The clue mentions a "nine{-}to{-}five routine," which is a common phrase referring to a standard workday. I will provide a five{-}letter word related to work or daily activities, such as "daily".} \\
\texttt{guess: daily} \\
            }
        }
    }
    & & \\ \\

    \theutterance \stepcounter{utterance}  
    & & \multicolumn{4}{p{0.6\linewidth}}{
        \cellcolor[rgb]{0.9,0.9,0.9}{
            \makecell[{{p{\linewidth}}}]{
                \texttt{\tiny{[GM$\rangle$P2]}}
                \texttt{You have attempted this same task before and did not succeed. These are the plans you wrote after those attempts.} \\
\texttt{Trial \#1: New plan: } \\
\\ 
\texttt{In my last attempt, I made a mistake by providing an agreement that was not based on the guess feedback. The guess feedback was not available at that point, and I should have stuck to the format and provided an agreement based on the clue and the guess. Next time, I will make sure to provide an agreement based on the available information, and if the guess feedback is not available, I will stick to the clue and the guess. For example, if the guess feedback is not available, I will provide an agreement based on whether the guess is directly related to the clue or not.} \\
\\ 
\texttt{Trial \#2: New plan: } \\
\\ 
\texttt{In my last attempt, I made a mistake by providing an agreement that was not based on the guess feedback. The guess feedback was not available at that point, and I should have stuck to the format and provided an agreement based on the clue and the guess. Next time, I will make sure to provide an agreement based on the available information, and if the guess feedback is not available, I will stick to the clue and the guess. For example, if the guess feedback is not available, I will provide an agreement based on whether the guess is directly related to the clue or not. However, I also need to consider the case when the guess is directly related to the clue but not necessarily correct. In such cases, I should not provide an agreement of "yes" without considering the guess feedback when it becomes available.} \\
\\ 
\\ 
\texttt{I need your assistance with a word game in which we need to find a 5{-}letter word using a clue, a guess and an explanation for the guess.} \\
\\ 
\texttt{For example, if the clue is "snow white" and the guess is "apple", you will get the information like this:} \\
\texttt{clue: snow white} \\
\texttt{explanation: the fairy tale Snow White involves the girl eating a poisoned apple, which has 5 letters.} \\
\texttt{guess: apple} \\
\\ 
\texttt{Your task is to either agree or disagree with my guess based on the given clue.} \\
\\ 
\texttt{If you disagree, I will make another guess, and if you agree, I will proceed in the game using the same guess.} \\
\\ 
\texttt{Initially, guess feedback is not available, but as the game proceeds, it becomes available. The guess feedback is not for current guess, but rather an indication of what letters can be present in the current guess. A letter highlighted in green indicates that it is a correct letter in the correct position, while yellow indicates that it is a correct letter in the wrong position, and red indicates a wrong letter. At this point, you should use all the provided information, including the clue, guess, explanation, and guess feedback, to decide whether the given guess matches the clue, and also stick to guess feedback or not.} \\
\\ 
\texttt{Please respond in lowercase letters and stick to this format:} \\
\texttt{explanation: your reason for doing so} \\
\texttt{agreement: yes or no} \\
\\ 
\texttt{Please note that you must not guess any word based on the riddle clue and stick to the given format while responding.} \\
\\ 
\texttt{Do not generate any other text.} \\
\\ 
\texttt{Let's begin!} \\
\\ 
\texttt{clue: nine{-}to{-}five routine} \\
\texttt{explanation: The clue mentions a "nine{-}to{-}five routine," which is a common phrase referring to a standard workday. I will provide a five{-}letter word related to work or daily activities, such as "daily".} \\
\texttt{guess: daily} \\
            }
        }
    }
    & \\ \\

    \theutterance \stepcounter{utterance}  
    & & & \multicolumn{4}{p{0.6\linewidth}}{
        \cellcolor[rgb]{0.9,0.9,0.9}{
            \makecell[{{p{\linewidth}}}]{
                \texttt{\tiny{[GM$\langle$P2]}}
                \texttt{explanation: the guess is directly related to the clue, which mentions a "nine{-}to{-}five routine" and the explanation mentions "daily activities". however, i should also consider the guess feedback when it becomes available, as the guess might be directly related to the clue but not necessarily correct.} \\
\\ 
\texttt{agreement: no} \\
            }
        }
    }
     \\ \\

    \theutterance \stepcounter{utterance}  
    & & \multicolumn{4}{p{0.6\linewidth}}{
        \cellcolor[rgb]{0.9,0.9,0.9}{
            \makecell[{{p{\linewidth}}}]{
                \texttt{\tiny{[P1$\langle$GM]}}
                \texttt{clue: nine{-}to{-}five routine} \\
\texttt{guess\_agreement: no} \\
\texttt{agreement\_explanation: the guess is directly related to the clue, which mentions a "nine{-}to{-}five routine" and the explanation mentions "daily activities". however, i should also consider the guess feedback when it becomes available, as the guess might be directly related to the clue but not necessarily correct.} \\
\\ 
\texttt{Provide your response only in this format.} \\
\texttt{explanation: details} \\
\texttt{guess: word} \\
            }
        }
    }
    & \\ \\

    \theutterance \stepcounter{utterance}  
    & \multicolumn{4}{p{0.6\linewidth}}{
        \cellcolor[rgb]{0.9,0.9,0.9}{
            \makecell[{{p{\linewidth}}}]{
                \texttt{\tiny{[P1$\rangle$GM]}}
                \texttt{explanation: The previous guess "daily" was not correct, but it was related to the clue. I will modify the guess by changing one letter at a time, as in the previous trial, to try to fit the clue. Since the previous guess was "daily", I will try to modify it by changing one letter at a time, for example, "dally".} \\
\texttt{guess: dally} \\
            }
        }
    }
    & & \\ \\

    \theutterance \stepcounter{utterance}  
    & & & \multicolumn{2}{p{0.3\linewidth}}{
        \cellcolor[rgb]{0.9,0.9,0.9}{
            \makecell[{{p{\linewidth}}}]{
                \texttt{\tiny{[GM$|$GM]}}
                \texttt{attempts = 1} \\
\texttt{target\_word = grind} \\
\texttt{target\_word\_clue = nine{-}to{-}five routine} \\
\texttt{guess = dally} \\
\texttt{guess\_feedback = d<yellow> a<red> l<red> l<red> y<red>} \\
            }
        }
    }
    & & \\ \\

    \theutterance \stepcounter{utterance}  
    & & \multicolumn{4}{p{0.6\linewidth}}{
        \cellcolor[rgb]{0.9,0.9,0.9}{
            \makecell[{{p{\linewidth}}}]{
                \texttt{\tiny{[P1$\langle$GM]}}
                \texttt{guess\_feedback: d<yellow> a<red> l<red> l<red> y<red>} \\
\\ 
\texttt{Provide your response only in this format.} \\
\texttt{explanation: details} \\
\texttt{guess: word} \\
            }
        }
    }
    & \\ \\

    \theutterance \stepcounter{utterance}  
    & \multicolumn{4}{p{0.6\linewidth}}{
        \cellcolor[rgb]{0.9,0.9,0.9}{
            \makecell[{{p{\linewidth}}}]{
                \texttt{\tiny{[P1$\rangle$GM]}}
                \texttt{explanation: The guess "dally" received feedback that the letter "d" is correct but in the wrong position, and the letters "a", "l", "l", and "y" are incorrect. I will try to modify the guess by keeping the correct letter "d" in a different position and replacing the incorrect letters with new ones that fit the clue. Since the previous guess was "dally" and the clue mentions a "nine{-}to{-}five routine", I will try to find a word related to work or daily activities that includes the letter "d".} \\
\texttt{guess: desk} \\
            }
        }
    }
    & & \\ \\

    \theutterance \stepcounter{utterance}  
    & & & \multicolumn{2}{p{0.3\linewidth}}{
        \cellcolor[rgb]{0.9,0.9,0.9}{
            \makecell[{{p{\linewidth}}}]{
                \texttt{\tiny{[GM$|$GM]}}
                \texttt{Error: The length of the guessed word is not 5.} \\
            }
        }
    }
    & & \\ \\

    \theutterance \stepcounter{utterance}  
    & & & \multicolumn{2}{p{0.3\linewidth}}{
        \cellcolor[rgb]{0.9,0.9,0.9}{
            \makecell[{{p{\linewidth}}}]{
                \texttt{\tiny{[GM$|$GM]}}
                \texttt{game\_result = ABORT} \\
            }
        }
    }
    & & \\ \\

    \theutterance \stepcounter{utterance}  
    & \multicolumn{4}{p{0.6\linewidth}}{
        \cellcolor[rgb]{0.95,0.95,0.95}{
            \makecell[{{p{\linewidth}}}]{
                \texttt{\tiny{[P2$\rightarrow$P1]}}
                \texttt{Post{-}episode Reflexion (not sent to GM):} \\
\texttt{New plan: The last episode got aborted because my guess "desk" is not a valid response. I should have produced a response in the format "explanation: details" followed by "guess: word". Moreover, in my last explanation I planned to modify the guess "dally" to fit the clue, but I didn't consider the fact that the previous guess "dally" received feedback that the letter "d" is correct but in the wrong position, and the letters "a", "l", "l", and "y" are incorrect. Next time, I will make sure to follow the correct format and provide a five{-}letter word related to work or daily activities, and also take into account the agreement with the previous guess and the guess feedback. Since the previous guess "dally" received feedback that the letter "d" is correct but in the wrong position, I will try to modify it by keeping the correct letter "d" in a different position and replacing the incorrect letters with new ones that fit the clue. For example, I could try to modify the guess "dally" to fit the clue, such as "dwell".} \\
            }
        }
    }
    & & \\ \\

    \theutterance \stepcounter{utterance}  
    & & & \multicolumn{4}{p{0.6\linewidth}}{
        \cellcolor[rgb]{0.95,0.95,0.95}{
            \makecell[{{p{\linewidth}}}]{
                \texttt{\tiny{[P1$\rightarrow$P2]}}
                \texttt{Post{-}episode Reflexion (not sent to GM):} \\
\texttt{New plan: } \\
\\ 
\texttt{In my last attempt, I made a mistake by providing an agreement that was not based on the guess feedback. The guess feedback was not available at that point, and I should have stuck to the format and provided an agreement based on the clue and the guess. Next time, I will make sure to provide an agreement based on the available information, and if the guess feedback is not available, I will stick to the clue and the guess. However, I also need to consider the case when the guess is directly related to the clue but not necessarily correct. In such cases, I should not provide an agreement of "yes" without considering the guess feedback when it becomes available. To improve my response, I will also consider the following: } \\
\\ 
\texttt{{-} If the guess is directly related to the clue, I will provide an agreement of "yes" only when the guess feedback is available and supports the guess.} \\
\texttt{{-} If the guess is not directly related to the clue, I will provide an agreement of "no" regardless of the guess feedback.} \\
\texttt{{-} If the guess is directly related to the clue but the guess feedback is not available, I will provide an agreement of "no" to avoid providing an incorrect agreement.} \\
            }
        }
    }
     \\ \\

\end{supertabular}
}

\end{document}
